\emph{Setting.} By Fact \claimref{pairing.imported-schwarzschild-connection-curvature} the static gauge has coframe $\theta^{0} = \sqrt f\,dt$, $\theta^{1} = dr/\sqrt f$, $\theta^{2} = r\,d\vartheta$, $\theta^{3} = r\sin\vartheta\,d\varphi$, so the physical frame acts on functions as $\hb(\gamma_{0}) = f^{-1/2}\partial_{t}$, $\hb(\gamma_{1}) = f^{1/2}\partial_{r}$, $\hb(\gamma_{2}) = r^{-1}\partial_{\vartheta}$, $\hb(\gamma_{3}) = (r\sin\vartheta)^{-1}\partial_{\varphi}$, and the rotation-gauge field along the frame, $\omega_{c} := \Omega(\hb(\gamma_{c}))$, read from the Fact's connection forms in the GTG sign of \claimref{obj.synthesis.gtg-flat-spacetime-gauge-fields}, is
\[
\omega_{0} = -\tfrac{f'}{2\sqrt f}\,\gamma_{0}\gamma_{1}, \qquad \omega_{1} = 0, \qquad \omega_{2} = \tfrac{\sqrt f}{r}\,\gamma_{1}\gamma_{2}, \qquad \omega_{3} = \tfrac{\sqrt f}{r}\,\gamma_{1}\gamma_{3} + \tfrac{\cot\vartheta}{r}\,\gamma_{2}\gamma_{3}.
\]
The check confirms that this field is torsion-free and has these values at its jet points. With $\gamma^{0} = \gamma_{0}$, $\gamma^{i} = -\gamma_{i}$, the covariant vector derivative is $D = \sum_{c}\gamma^{c} D_{\hb(\gamma_{c})}$, with $D_{a}\psi = a \cdot \nabla\psi + \tfrac12\Omega(a)\psi$ on spinors and $D_{a}M = a \cdot \nabla M + \Omega(a) \times M$ on multivectors \claimref{obj.synthesis.gtg-flat-spacetime-gauge-fields}. Write $\partial_{r_{*}} = f\partial_{r}$.

\emph{(2), the Dirac operator.} The spinor connection term is $\Gamma := \tfrac12\sum_{c}\gamma^{c}\omega_{c}$. Term by term: $\tfrac12\gamma^{0}\omega_{0} = -\tfrac{f'}{4\sqrt f}\gamma_{1} = \tfrac{f'}{4\sqrt f}\gamma^{1}$; $\tfrac12\gamma^{2}\omega_{2} = -\tfrac{\sqrt f}{2r}\gamma_{2}\gamma_{1}\gamma_{2} = \tfrac{\sqrt f}{2r}\gamma^{1}$; and $\tfrac12\gamma^{3}\omega_{3} = \tfrac{\sqrt f}{2r}\gamma^{1} + \tfrac{\cot\vartheta}{2r}\gamma^{2}$. So $\Gamma = \gamma^{1}\bigl(\tfrac{f'}{4\sqrt f} + \tfrac{\sqrt f}{r}\bigr) + \gamma^{2}\tfrac{\cot\vartheta}{2r}$, a vector (no trivector part), and
\[
D\psi = \gamma^{0}f^{-1/2}\partial_{t}\psi + \gamma^{1}\Bigl(f^{1/2}\partial_{r} + \tfrac{f'}{4 f^{1/2}} + \tfrac{f^{1/2}}{r}\Bigr)\psi + \tfrac1r\Bigl(\gamma^{2}\bigl(\partial_{\vartheta} + \tfrac12\cot\vartheta\bigr) + \gamma^{3}\tfrac{1}{\sin\vartheta}\partial_{\varphi}\Bigr)\psi .
\]
For $\psi = r^{-1}f^{-1/4}\Phi$, $\partial_{r}(r^{-1}f^{-1/4}) = -r^{-1}f^{-1/4}\bigl(\tfrac1r + \tfrac{f'}{4f}\bigr)$ cancels the two zero-order terms of the radial bracket, so $\bigl(f^{1/2}\partial_{r} + \tfrac{f'}{4f^{1/2}} + \tfrac{f^{1/2}}{r}\bigr)\psi = r^{-1}f^{-1/4}f^{1/2}\partial_{r}\Phi$, and multiplying by $r f^{3/4}$,
\[
r f^{3/4}\,D\psi = \gamma^{0}\partial_{t}\Phi + \gamma^{1}\partial_{r_{*}}\Phi + \tfrac{\sqrt f}{r}\,K\Phi = A\Phi + WK\Phi .
\]
The Dirac-Hestenes equation $D\psi\,I\sigma_{3} = m\psi\gamma_{0}$ ($I\sigma_{3} = \gamma_{2}\gamma_{1}$) becomes $(A\Phi + WK\Phi)I\sigma_{3} = m\sqrt f\,\Phi\gamma_{0}$, since $r f^{3/4}\cdot r^{-1}f^{-1/4} = \sqrt f$. The check verifies the displayed reduction for a generic even multivector field (its second-order jet) at two points, and that the wrong measure factor $r^{-1}f^{+1/4}$ fails.

\emph{Channels.} Let $\gamma^{01} := \gamma^{0}\gamma^{1}$. It anticommutes with $\gamma^{0}$, $\gamma^{1}$, commutes with $\gamma^{2}$, $\gamma^{3}$, and $(\gamma^{01})^{2} = 1$. $A$ involves only $\partial_{t}, \partial_{r}$ and the vectors $\gamma^{0}, \gamma^{1}$; $K$ only $\partial_{\vartheta}, \partial_{\varphi}$ with $\vartheta$-dependent coefficients and $\gamma^{2}, \gamma^{3}$; so the differential parts commute and the Clifford parts anticommute, $KA = -AK$, and $\gamma^{01}KA = -\gamma^{01}AK = A\gamma^{01}K$. Both $A$ and $K$ are left multiplications by vectors, so each maps even fields to odd ones; $\gamma^{01}K$ alone has no even eigenvectors. Define the sphere operator on even fields by a right factor $\gamma_{0}$:
\[
\tilde T\Phi := \gamma^{01}(K\Phi)\,\gamma_{0} .
\]
It maps even fields to even fields, commutes with $A$ (which acts on the left), with right multiplication by $I\sigma_{3} = \gamma_{2}\gamma_{1}$ and by $\gamma_{0}$ (which commute with each other), with $\partial_{t}$ and $\partial_{\varphi}$, and $\tilde T^{2}\Phi = \gamma^{01}K\gamma^{01}K\Phi\gamma_{0}\gamma_{0} = K^{2}\Phi$. The eigenfunctions of $\tilde T$ are taken to span the even multivector fields on the sphere: this is the classical input (C1) of the formal level (from memory: $\tilde T$ is a first-order elliptic operator on $S^{2}$, symmetric for the $L^{2}$ pairing with measure $\sin\vartheta\,d\vartheta\,d\varphi$, with the spinor spherical harmonics as eigenfunctions and eigenvalues $\pm(j + \tfrac12)$; only the decomposition is used). On an eigenspace $\tilde T\Phi = \lambda\Phi$, multiplying by $\gamma^{01}$ on the left and $\gamma_{0}$ on the right gives $K\Phi = \lambda\gamma^{01}\Phi\gamma_{0}$, which turns the Dirac-Hestenes equation of (2) into the channel equation
\[
\bigl(A\Phi + \lambda W\gamma^{01}\Phi\gamma_{0}\bigr)I\sigma_{3} = m\sqrt f\,\Phi\gamma_{0},
\]
every term of which is even. For $m = 0$, with $C\Phi := A\Phi + \lambda W\gamma^{01}\Phi\gamma_{0}$,
\[
C(C\Phi) = A^{2}\Phi + \lambda(\partial_{r_{*}}W)\gamma^{1}\gamma^{01}\Phi\gamma_{0} - \lambda W\gamma^{01}(A\Phi)\gamma_{0} + \lambda W\gamma^{01}(A\Phi)\gamma_{0} + \lambda^{2}W^{2}\Phi ,
\]
with $A^{2} = \partial_{t}^{2} - \partial_{r_{*}}^{2}$ and $\gamma^{1}\gamma^{01} = \gamma^{0}$; so $(\partial_{t}^{2} - \partial_{r_{*}}^{2} + \lambda^{2}W^{2})\Phi + \lambda(\partial_{r_{*}}W)\,\gamma^{0}\Phi\gamma_{0} = 0$. The map $\Phi \mapsto \gamma^{0}\Phi\gamma_{0}$ is an involution of the even fields, commuting with $A^{2}$, and on its $\pm1$ eigenspaces the potentials are $V_{\pm} = \lambda^{2}W^{2} \pm \lambda\,\partial_{r_{*}}W$. Since $\partial_{r_{*}}(\sqrt f/r) = f\bigl(\tfrac{f'}{2\sqrt f\,r} - \tfrac{\sqrt f}{r^{2}}\bigr) = \sqrt f\bigl(\tfrac{f'}{2r} - \tfrac{f}{r^{2}}\bigr) = \sqrt f\,\tfrac{3M - r}{r^{3}}$, $V_{\pm} = \lambda^{2}f/r^{2} \pm \lambda\sqrt f(3M - r)/r^{3}$. The check verifies on a generic third-order jet that $\tilde T\Phi$ is even, that $\tilde T$ commutes with $A$ and with the two right multiplications, that $\tilde T^{2} = K^{2}$, and the squared channel identity for $\lambda = 1, -2, 3$; it also verifies that the identity fails without the right factor $\gamma_{0}$.

\emph{(1), the scalar.} For a function $\phi$, $D\phi = \sum_{a}\gamma^{a}\,\hb(\gamma_{a})\phi$, and for a vector field $V$, $D \cdot V = \sum_{c}\gamma^{c} \cdot (\hb(\gamma_{c})V + \omega_{c} \times V)$; since $\Omega$ is torsion-free, this is the divergence of the derived metric, $D \cdot V = (-g)^{-1/2}\partial_{\mu}\bigl((-g)^{1/2}V^{\mu}\bigr)$ with $(-g)^{1/2} = r^{2}\sin\vartheta$ (the connection terms are the divergence of the frame). Collecting, $D \cdot D\phi = f^{-1}\partial_{t}^{2}\phi - r^{-2}\partial_{r}(r^{2}f\partial_{r}\phi) - r^{-2}\Delta_{S^{2}}\phi$, the d'Alembertian of the derived metric in the signature $(+,-,-,-)$. With $\phi = r^{-1}u\,Y_{\ell m}$, $\Delta_{S^{2}}Y_{\ell m} = -\ell(\ell+1)Y_{\ell m}$ and $r^{-2}\partial_{r}(r^{2}f\partial_{r}(u/r)) = r^{-1}\partial_{r}(f\partial_{r}u) - f'u/r^{2}$,
\[
f\,r\,(D \cdot D\phi + \mu^{2}\phi) = \Bigl(\partial_{t}^{2}u - f\partial_{r}(f\partial_{r}u) + f\bigl(\tfrac{\ell(\ell+1)}{r^{2}} + \tfrac{f'}{r} + \mu^{2}\bigr)u\Bigr)Y_{\ell m},
\]
and $f'/r = 2M/r^{3}$. The check computes $D \cdot D\phi$ in GA from the field of the Fact (gradient, then covariant divergence with $\omega_{c}$) at a jet point, for $\ell = 1$ ($Y = \cos\vartheta$) and $\ell = 2$ ($Y = \sin\vartheta\cos\vartheta\cos\varphi$), with $\mu^{2} = \tfrac19$, and compares with the right side.

\emph{(3), Maxwell.} One-forms on the sphere split into exact and co-exact parts (the classical input (C2) of the formal level); the co-exact (odd-parity, axial) part of a one-form of angular type $(\ell, m)$, $\ell \geq 1$, is $\psi\,X_{\ell m}$ with $X_{\ell m} = (\sin\vartheta)^{-1}\partial_{\varphi}Y_{\ell m}\,d\vartheta - \sin\vartheta\,\partial_{\vartheta}Y_{\ell m}\,d\varphi$, and gauge transformations $d\chi$ are exact, so the odd-parity sector is described by the single function $\psi(t,r)$ with no gauge freedom; the $dt$ and $dr$ components of an odd-parity potential vanish, since scalar harmonics are polar. Read $A_{\ell m}$ as the physical vector $\sum_{\mu}(A_{\ell m})_{\mu}g^{\mu}$, $g^{\mu} = \hb^{T}(e^{\mu})$, and set $F = D \wedge A_{\ell m}$; since $\Omega$ is torsion-free, $F$ is the exterior derivative of $A_{\ell m}$ and $D \wedge F = 0$. In components, $F_{tA} = \partial_{t}\psi\,(X_{\ell m})_{A}$, $F_{rA} = \partial_{r}\psi\,(X_{\ell m})_{A}$ ($A \in \{\vartheta, \varphi\}$) and $F_{\vartheta\varphi} = \psi\,(\partial_{\vartheta}X_{\varphi} - \partial_{\varphi}X_{\vartheta}) = -\psi\sin\vartheta\,\Delta_{S^{2}}Y_{\ell m} = \ell(\ell+1)\psi\sin\vartheta\,Y_{\ell m}$. The divergence $D \cdot F$ has components $(-g)^{-1/2}\partial_{\mu}((-g)^{1/2}F^{\mu\nu})$, with $g^{tt} = 1/f$, $g^{rr} = -f$, $g^{\vartheta\vartheta} = -r^{-2}$, $g^{\varphi\varphi} = -(r\sin\vartheta)^{-2}$. For $\nu = t, r$ it is a multiple of the sphere divergence of $X_{\ell m}$, which vanishes (a co-exact form is co-closed); so does $D \cdot A_{\ell m}$. For $\nu = \varphi$, the three terms of $\partial_{\mu}((-g)^{1/2}F^{\mu\varphi})$ are $-\tfrac{1}{f\sin\vartheta}\partial_{t}^{2}\psi\,X_{\varphi}$, $\tfrac{1}{\sin\vartheta}\partial_{r}(f\partial_{r}\psi)\,X_{\varphi}$ and $\tfrac{1}{r^{2}}\partial_{\vartheta}\bigl(F_{\vartheta\varphi}/\sin\vartheta\bigr) = \tfrac{\ell(\ell+1)}{r^{2}}\psi\,\partial_{\vartheta}Y_{\ell m} = -\tfrac{\ell(\ell+1)}{r^{2}\sin\vartheta}\psi\,X_{\varphi}$; multiplied by $-f\sin\vartheta$ they give $\bigl(\partial_{t}^{2}\psi - f\partial_{r}(f\partial_{r}\psi) + f\ell(\ell+1)\psi/r^{2}\bigr)X_{\varphi}$, and $\nu = \vartheta$ is the same with $X_{\vartheta}$. In GA,
\[
D \cdot F = \frac{1}{f}\Bigl(\partial_{t}^{2}\psi - f\partial_{r}(f\partial_{r}\psi) + f\,\frac{\ell(\ell+1)}{r^{2}}\psi\Bigr)\frac{A_{\ell m}}{\psi},
\]
the odd-parity Maxwell channel equation. The check computes $F = D \wedge A$, $D \wedge F$, $D \cdot A$ and $D \cdot F$ in GA from the field of the Fact at a jet point, for $\ell = 1$ and $\ell = 2$ and two second-order jets of $\psi$, and finds this identity exactly (the identity is linear in $\psi$ and of second order, so a generic second-order jet tests it at that point for every $\psi$). For the even parity: the pseudoscalar $I$ is constant and commutes with every bivector, so $D_{a}(IF) = I D_{a}F$, and $\gamma^{c}I = -I\gamma^{c}$, so $D(IF) = -I\,DF$ (checked). Duality $F \mapsto IF$ is invertible on source-free fields ($I^{2} = -1$) and exchanges the two parity sectors ($I$ is odd under reflection), so every even-parity source-free field of type $(\ell, m)$ is $IF$ for an odd-parity one, and its channel equation is the same equation for $\psi$.

\emph{(4), at the horizon.} With $r_{*} = \int dr/f$, $\tfrac{d}{dr}(\ln f - 2\kappa r_{*}) = (f' - 2\kappa)/f$, and with $\kappa = f'(2M)/2 = 1/(4M)$ this is $\bigl(\tfrac{2M}{r^{2}} - \tfrac{1}{2M}\bigr)\tfrac{r}{r - 2M} = -\tfrac{r + 2M}{2Mr}$, a rational function regular at $r = 2M$ (the check computes it for $M = 1$: $-(r+2)/(2r)$, value $-1$ at the horizon). So $\ln f - 2\kappa r_{*}$ has a finite limit as $r \to 2M$, and $f = C(r)e^{2\kappa r_{*}}$ with $C$ regular and positive there. The scalar potential is $f$ times $\ell(\ell+1)/r^{2} + 2M/r^{3} + \mu^{2}$, which is $(\ell(\ell+1) + 1)/(4M^{2}) + \mu^{2} > 0$ at $r = 2M$; the Maxwell potential is $f$ times $\ell(\ell+1)/r^{2}$; both decay as $e^{2\kappa r_{*}}$ as $r_{*} \to -\infty$. The Dirac couplings $\lambda W = \lambda\sqrt f/r$ and $m\sqrt f$ decay as $e^{\kappa r_{*}}$, and $V_{\pm}/\sqrt f = \lambda^{2}\sqrt f/r^{2} \pm \lambda(3M - r)/r^{3} \to \pm\lambda/(8M^{2})$: exact order $\sqrt f$. Hence for each fixed channel the coefficients of the channel equation that couple it to the sphere, to the mass and to the curvature tend to zero as $r_{*} \to -\infty$, and the channel equation tends to the free two-dimensional equation: $(\partial_{t}^{2} - \partial_{r_{*}}^{2})u = 0$, solved by $u_{+}(t - r_{*}) + u_{-}(t + r_{*})$, and $A\Phi = 0$, solved by even $\Phi(t - r_{*})$ with $(\gamma^{0} - \gamma^{1})\Phi' = 0$, that is $\gamma^{01}\Phi = \Phi$ (left multiplication by $\gamma^{01}$ preserves even fields) (the idempotent $\tfrac12(1 + \gamma^{01})$ is annihilated by $\gamma^{0} - \gamma^{1}$, checked), and by $\Phi(t + r_{*})$ with $\gamma^{01}\Phi = -\Phi$. Each channel is therefore a two-dimensional mode at the horizon: this is the channel-reduction clause of joint J3; the identification of the s-wave theory's current boson with these channels, and spin 2, are not addressed. \emph{Scope (the frame).} The statement is made in the static frame. With $u_{\mathrm{s}} = t - r_{*}$ and $v_{\mathrm{s}} = t + r_{*}$, $\partial_{t}^{2} - \partial_{r_{*}}^{2} = 4\partial_{u_{\mathrm{s}}}\partial_{v_{\mathrm{s}}}$, and with $U_{\mathrm{K}} = -e^{-\kappa u_{\mathrm{s}}}$, $V_{\mathrm{K}} = e^{\kappa v_{\mathrm{s}}}$ one has $\partial_{u_{\mathrm{s}}}\partial_{v_{\mathrm{s}}} = -\kappa^{2}U_{\mathrm{K}}V_{\mathrm{K}}\,\partial_{U_{\mathrm{K}}}\partial_{V_{\mathrm{K}}} = \kappa^{2}e^{2\kappa r_{*}}\,\partial_{U_{\mathrm{K}}}\partial_{V_{\mathrm{K}}}$. So a channel equation $(\partial_{t}^{2} - \partial_{r_{*}}^{2} + V)u = 0$ with $V = f\,v(r)$ becomes $\partial_{U_{\mathrm{K}}}\partial_{V_{\mathrm{K}}}u = -(4\kappa^{2})^{-1}e^{-2\kappa r_{*}}f\,v\,u = -(4\kappa^{2})^{-1}C(r)\,v(r)\,u$, with coefficient $-4M^{2}C(2M)v(2M)$ at the horizon: nonzero for the scalar channel ($v(2M) = (\ell(\ell+1) + 1)/(4M^{2}) + \mu^{2}$) and for the Maxwell channel ($v(2M) = \ell(\ell+1)/(4M^{2})$, $\ell \geq 1$). For $r_{*} = r + 2M\ln(r/2M - 1)$, $e^{2\kappa r_{*}} = (r/2M - 1)e^{r/2M}$ and $f = (2M/r)(r/2M - 1)$, so $C(r) = (2M/r)e^{-r/2M}$ and $C(2M) = e^{-1}$. The decay $e^{2\kappa r_{*}}$ of the potential is therefore the statement that $V/f$ is regular at the horizon, read in Killing null coordinates; it is not a decoupling in coordinates regular across the horizon. This paragraph is a direct substitution and is not part of the check; for the Dirac channel only the static-frame statement is made. The limit is not uniform in $\ell$: at fixed $r$ the potentials grow like $\ell(\ell+1)$ or $\lambda^{2}$, so the region in which a channel is free recedes towards the horizon as the angular momentum grows (argued).

\emph{Remark (prior art, from memory; not used).} The potentials are those of Regge and Wheeler (spin $0$ and odd-parity spin $1$; the even parity of spin $1$ has the same potential) and Chandrasekhar's $W^{2} \pm dW/dr_{*}$ for the Dirac field; Lasenby, Doran and Gull separated the Dirac equation in the Schwarzschild field in GTG in the Painlev\'e-Gullstrand gauge. In the Teukolsky form, with a tetrad singular at the future horizon, the spin-weighted radial potentials do not vanish there; J3 concerns channels rescaled to boost weight zero, as here, where the factor $f^{-1/4}$ of the Dirac channel is that rescaling for weight $\tfrac12$ (argued).
